% !TEX root = ../main.tex
% ---------------------------------------------------------------------
%  P6 -- Rekonfigurierbare Hardware
% ---------------------------------------------------------------------
\vlsec{cP6}{P6}{Rekonfigurierbare Hardware}
\label{sec:reconfig}

\SW{Spektrum} fein\,\ra\,grob, schnell\,\ra\,flexibel:
\begin{ztab}{@{}lXl@{}}
Ebene & Baustein & Eigenschaft \\ \midrule
Gatter & \T{FPGA} & sehr feingranular \\
Sequencer & progr.\ FSM & kaum progr., sehr schnell \\
 & Controller & winziger Befehlssatz (Load, Store, Add,
   Jump, Compare), flexibel, langsam \\
 & \T{PIO} & flexibel \emph{und} schnell \\
 & kleine CPU $+$ RAM & sehr flexibel, sehr langsam \\
Operation & \T{CGRA} & routbare ALU/FPU-Blöcke, Datenpfad, VLIW \\
\end{ztab}

\subsection{FPGA}
\begin{itemize}
\item \T{LUT} (Look Up Table) realisiert \emph{jede} Logikfunktion ihrer
      Eingänge; im Beispiel 3, real 4--8
\item Konfigurationsspeicher $=$ SRAM-Zellen in Schieberegister-Anordnung
      (Wiederverwendung der Scan-Test-Struktur)
\item Switch-Matrix-Routing mit Sechs-Transistor-Kreuzungen
\item konfiguriert in VHDL
\end{itemize}
\begin{formel}[Die Zahl, die man sich merkt]
\T{Routing-Delay $>$70\,\% des Gesamt-Delays.} Nicht die Logik ist
langsam, sondern die Verdrahtung.
\end{formel}

\subsection{CGRA}
Array aus \T{PE} (Processing Elements) mit \T{FU} (spezialisierte
ALU/FPU: AND, SHL, ADD, MUL, MAC), typisch extremes \T{VLIW}.

\begin{bsp}[4-Tap-FIR auf $4\times4$-CGRA]
\label{sec:cgra}
$y[n]=h_0x[n]+h_1x[n{-}1]+h_2x[n{-}2]+h_3x[n{-}3]$,
Nearest-Neighbor-Mesh, 1-Takt-ADD/MUL.
\[m_k=h_k\,x[n{-}k],\quad a_0=m_0{+}m_1,\ a_1=m_2{+}m_3,\ y=a_0{+}a_1\]
\T{7 von 16 PE} benutzt $\;\;$ \T{1 Ergebnis pro Takt} (Pipeline)
$\;\;$ Latenz $=1_{\text{MUL}}+2_{\text{ADD}}=\T{3 Takte}$.
\end{bsp}

\SW{Xilinx AI Engines} Kachel aus Vektorprozessor (SIMD für parallele
Arithmetik) $+$ Skalarprozessor (Kontrollfluss, Adressen) $+$ lokalem
deterministischem Speicher $+$ DMA-Engines, verbunden über streaming-
und memory-mapped Interconnect.

\subsection{RP2350 PIO}
3 identische PIO-Blöcke, je \T{4 State Machines}. Realisiert in Software
Schnittstellen wie I2C, 3-pin I2S, SDIO, SPI/DSPI/QSPI, UART, DPI,
8080/6800-Parallelbus, VGA (Widerstands-DAC).
\begin{ztab}{@{}lX@{}}
Register je SM & $2\times32$\,Bit Shift (\T{OSR}, \T{ISR}),
  $2\times32$\,Bit Scratch (\T{X}, \T{Y}) \\
FIFO & $4\times32$\,Bit je TX und RX \\
Instruktionsspeicher & 32 Worte, gemeinsam für die 4 SM \\
Takt & fraktionaler Teiler (16 Integer-, 8 Nachkommabits) \\
Rest & flexibles GPIO-Mapping, DMA-Interface, IRQ-Flags \\
\end{ztab}

\begin{formel}[Nur 9 Befehle]
\R{JMP} \R{WAIT} \R{IN} \R{OUT} \R{PUSH} \R{PULL} \R{MOV} \R{IRQ}
\R{SET}. \R{NOP} ist eine Pseudo-Instruktion ($=$ \R{mov y,y}).\\
Format: \R{<instr> (side <wert>) ([<delay>])}, Delay 0--31 Takte
(weniger, wenn \R{.side\_set} benutzt wird).\\
JMP-Bedingungen: immer, \R{!X}, \R{X--}, \R{!Y}, \R{Y--}, \R{X!=Y},
\R{PIN}, \R{!OSRE}.
\end{formel}

\SW{Datenfluss} \R{PULL} holt ein Wort aus der TX-FIFO in den \T{OSR},
\R{OUT} schiebt 1--32\,Bit daraus zum Ziel; \R{IN} schiebt in den
\T{ISR}, \R{PUSH} schreibt ihn in die RX-FIFO.
\T{Autopull/Autopush} füllt bzw.\ leert automatisch
\ra die \R{pull}/\R{push}-Instruktion entfällt.

\SW{Zwei Tricks, die Takte sparen}
\begin{description}
\item[Side-Set] schaltet \emph{parallel} zum Befehl Pins um
  \ra Daten \emph{und} Takt ohne Zusatzbefehl (SPI).
\item[Program Wrapping] \R{.wrap\_target}/\R{.wrap} ersetzt den \R{jmp}
  durch einen Rücksprung, der \emph{null} Takte kostet.
\end{description}

\begin{lstlisting}
.program squarewave_wrap
  set pindirs, 1     ; Pin auf Ausgang
.wrap_target         ; 0-Takt-Ruecksprungziel
  set pins, 1 [1]    ; high, 1 Takt warten
  set pins, 0 [1]    ; low,  1 Takt warten
.wrap                ; 0-Takt-Sprung statt jmp
\end{lstlisting}

\SW{Direktiven} \R{.program}, \R{.define}, \R{.side\_set <n> (opt)},
\R{.in}/\R{.out <n> (left|right) (auto) (<schwelle>)},
\R{.fifo}, \R{.clock\_div}, \R{.mov\_status},
\R{.wrap\_target}/\R{.wrap}, \R{.origin}, \R{.pio\_version}
(1 $=$ RP2350, 0 $=$ RP2040), \R{.word}.
